\documentclass[../../../include/open-logic-section]{subfiles}

\begin{document}

\olfileid{sth}{ordinals}{vn}	
\olsection[వాన్ న్యూమన్ నిర్మాణం]{వాన్ న్యూమన్ క్రమసంఖ్యల నిర్మాణం}

\olref[sth][ordinals][iso]{thm:woalwayscomparable}
ఒక ఆలోచనకు దారితీస్తుంది.
సుక్రమాలకు ప్రతినిధులుగా
\emph{క్రమరకాలు} అనే
కొన్ని వస్తువులను
పరిచయం చేయవచ్చు.
$\tuple{A, <}$ సుక్రమపు
క్రమరకాన్ని $\ordtype{A, <}$
అని రాస్తే, ఈ రెండు
సూత్రాలు లభించాలని
ఆశిస్తాం:
\begin{align*}
	\ordtype{A, <} = \ordtype{B, \lessdot} & 
	\text{ అప్పుడే మరియు అప్పుడే } \ordeq{\tuple{A, <}}{\tuple{B, \lessdot}}\\
	\ordtype{A, <} < \ordtype{B, \lessdot}&
	\text{ అప్పుడే మరియు అప్పుడే }\ordeq{\tuple{A, <}}{\tuple{B_b, \lessdot_b}}\text{ ఏదో ఒక }b \in B
\end{align*}
ఇంకా, సహజ సంఖ్యలను
కొన్ని సమితులుగా
పరిచయం చేసినట్టే,
క్రమరకాలను కూడా
\emph{కొన్ని సమితులుగా}
పరిచయం చేయగలమని
ఆశించవచ్చు.

దీనికి అత్యంత సాధారణ
పద్ధతి---మనం అనుసరించబోయే
పద్ధతి---కొన్ని \emph{నియత}
సుక్రమిత సమితుల ద్వారా
ఈ క్రమరకాలను
నిర్వచించడం. ఈ నియత
సమితులను మొదట
వాన్ న్యూమన్ పరిచయం
చేశాడు:

\begin{defn}
$(\forall x \in A)x \subseteq
A$ అయితేనే $A$ సమితి
\emph{సంక్రామకం}.
$A$ సంక్రామకమై,
$\in$ ద్వారా సుక్రమితమైతేనే
$A$ ఒక \emph{క్రమసంఖ్య}.
\end{defn}
\noindent
ఇకపై క్రమసంఖ్యలకు
గ్రీకు అక్షరాలను వాడతాం.
నిర్వచనం నుంచి వెంటనే:
$\alpha$ క్రమసంఖ్య అయితే,
$\tuple{\alpha, \in_\alpha}$
ఒక సుక్రమం; ఇక్కడ
$\in_\alpha =
\Setabs{\tuple{x, y} \in \alpha^2}{x \in y}$.
కాబట్టి సంకేతలిపిని
కొంచెం విస్తృతార్థంలో
వాడుతూ $\alpha$
\emph{స్వయంగా} సుక్రమం
అని చెప్పవచ్చు.

మన తొలి కొన్ని
క్రమసంఖ్యలు ఇవి:
\[
	\emptyset, \{\emptyset\}, 
	\{\emptyset, \{\emptyset\}\}, 
	\{\emptyset, \{\emptyset\}, \{\emptyset, \{\emptyset\}\}\}, \ldots
\]
ఇవే అనంతత్వ స్వీకృతంలో,
అంటే $\omega$కు వాన్
న్యూమన్ ఇచ్చిన నిర్వచనంలో,
మనం చూసిన తొలి
క్రమసంఖ్యలని గమనించండి
(\olref[sth][z][infinity-again]{sec}
చూడండి). ఇది యాదృచ్ఛికం
కాదు. వాన్ న్యూమన్
క్రమసంఖ్యల నిర్వచనం
సహజ సంఖ్యలను
క్రమసంఖ్యలుగా పరిగణిస్తుంది;
అయితే అతిపరిమిత
క్రమసంఖ్యలనూ అనుమతిస్తుంది.

ఎప్పటిలాగే ఇప్పుడు
ఇలా అడగవచ్చు:
\emph{ఇవేనా} క్రమసంఖ్యలు?
లేక వాన్ న్యూమన్
క్రమసంఖ్యలు\emph{గా
పరిగణించగల} కొన్ని
సమితులను మాత్రమే
ఇచ్చాడా? ఈ ప్రశ్నపై
చర్చ \olref[sfr][rel][ref]{sec},
\olref[sfr][arith][ref]{sec},
\olref[sfr][infinite][dedekindsproof]{sec},
\olref[sth][z][nat]{sec}లలోని
చర్చల మాదిరిగానే
ఉంటుంది; కాబట్టి
దీనిపై ఎక్కువగా
ఆగం. ఇకపై
``క్రమసంఖ్యలు'' అంటే
``వాన్ న్యూమన్
క్రమసంఖ్యలు'' అనే
అర్థంలోనే వాడతాం.

\end{document}
